\begin{question}
\noindent The diagram shows three straight lines meeting at point $C$. $A$, $C$ and $B$ lie on a straight line.

\begin{center}
\begin{tikzpicture}[scale=1]
    % Define the points
    \coordinate (A) at (-3,0);
    \coordinate (B) at (3,0);
    \coordinate (C) at (0,0);
    \coordinate (D) at ($(C)+(65:2cm)$);
    \coordinate (E) at ($(C)+(140:2cm)$);

    % Draw the straight line AB
    \draw (A) -- (B);
    % Draw the two rays from C
    \draw (C) -- (D);
    \draw (C) -- (E);

    % Angle markings
    \pic [draw, "$65^{\circ}$", angle eccentricity=1.6, angle radius=0.7cm] {angle = B--C--D};
    \pic [draw, "$x^{\circ}$", angle eccentricity=1.7, angle radius=0.8cm] {angle = D--C--E};
    \pic [draw, "$40^{\circ}$", angle eccentricity=1.8, angle radius=0.6cm] {angle = E--C--A};

    % Labels
    \node at (A) [left,font=\small] {$A$};
    \node at (B) [right,font=\small] {$B$};
    \node at (C) [below,font=\small] {$C$};
    \node at (D) [above,font=\small] {$D$};
    \node at (E) [above,font=\small] {$E$};
\end{tikzpicture}
\end{center}

\noindent Work out the size of angle $x$.
\vspace{4cm}
\begin{flushright}
\answerequnit{x}{$^\circ$}{2}
\end{flushright}
\end{question}